\section{总结与延伸}

Model-Based RL 是 RL 工具箱中的重要组成部分，其核心优势在于数据效率。

关键要点：
\begin{enumerate}
    \item 学习动力学模型可以大幅减少所需的真实交互数据
    \item 模型可用于 planning（CEM/MPC）或生成合成数据（Dyna 风格）
    \item 模型误差是核心挑战，需要通过集成、短视野 rollout、持续更新来缓解
    \item 在机器人操作等数据昂贵的场景中特别有价值
\end{enumerate}

\begin{knowledgebox}{拓展阅读}
\begin{itemize}
    \item Nagabandi et al., ``Neural Network Dynamics for Model-Based Deep RL with Model-Free Fine-Tuning,'' ICRA 2018
    \item Chua et al., ``Deep Reinforcement Learning in a Handful of Trials using Probabilistic Dynamics Models (PETS),'' NeurIPS 2018
    \item Janner et al., ``When to Trust Your Model: Model-Based Policy Optimization (MBPO),'' NeurIPS 2019
    \item Sutton, ``Dyna, an Integrated Architecture for Learning, Planning, and Reacting,'' 1991
\end{itemize}
\end{knowledgebox}

\end{document}
